\documentclass[11pt]{article}

\usepackage[margin=1in]{geometry}
\usepackage{amsmath}
\usepackage{amssymb}
\usepackage{hyperref}
\usepackage{enumitem}
\usepackage{booktabs}

\setlength{\parindent}{0pt}
\setlength{\parskip}{0.5em}

\newcommand{\hccultra}{\textsc{HccUltraFit}}

\title{Two undergraduate projects:\\ stratified sampling for redistricting plans}
\author{}
\date{\today}

\begin{document}
\maketitle

\section{Context}

These projects build on \emph{Towards stratified sampling for
redistricting plans} (Wang, Herschlag, Yim, Mattingly, Gilbert), which
constructs a stratification of the space of redistricting plans:
districts pooled from an ensemble are clustered into representative
``letters'' with \hccultra{}~\cite{yim2024}, plans are encoded as
unordered tuples of letters (``words''), and a partition of unity over
a representative word set yields stratum weights $z_s$ and a flux
matrix $F$. The paper stops short of two things: actually running
stratified (umbrella) sampling on these strata, and connecting the
strata to the election outcomes that ensembles are used to study in
litigation~\cite{herschlag2020}. Each project takes one of these gaps.

Both projects work at two scales.

\textbf{Small scale (exact): the Duplin--Onslow system.} A 25-unit
dual graph (24 Onslow County, NC precincts plus Duplin County as a
single unit) from the Duke Quantifying Gerrymandering project, with
real population, Black voting-age population, and precinct returns for
22 races (2008--2016) attached, together with an exhaustive
enumeration of \emph{all} 17{,}653 population-balanced contiguous
3-district plans. Because the plan space is fully enumerated, every
pipeline quantity --- letters, words, POU, $z_s$, $F$, and every
election observable --- is exact. There is no sampling error anywhere:
every claim is checkable against ground truth. The full pipeline
already runs on this system in the \texttt{gmstrat} repository (see
Section~\ref{sec:infra}).

\textbf{Full scale (sampled): Connecticut and North Carolina.} The
real-state pipelines from the paper ($|V|=462$, $I=5$ and $|V|=2671$,
$I=14$), sampled with Metropolized Cycle Walk~\cite{cyclewalk}, plus
the $m\times n$ grid-graph ladder in \texttt{gmstrat}
(\texttt{data/graph/}, $6\times6$ through $20\times20$) for controlled
scaling experiments between the two regimes.

\subsection*{Established starting points (computed June 2026, exact)}

On Duplin--Onslow, with the population-weighted symmetric-difference
metric, $K=12$ letters, beam width $L=10$ (49 strata):

\begin{itemize}[itemsep=2pt]
  \item The 14{,}345 distinct districts fill district space densely:
    within-cluster dispersion only halves by $K=80$, with the
    interpretable structure exhausted by $K\approx 10$--$12$. Letters
    at $K=12$ are geographically legible motifs (a Duplin-anchored
    type, Jacksonville-centered, coastal-south, northern-tier).
  \item The strata communicate well: exact flux matrix with spectral
    gap $1-|\lambda_2| \approx 0.48$; largest exact stratum weight
    $z_s \approx 0.074$.
  \item \textbf{Alignment with election outcomes is weak.} The
    between-stratum share of variance (an alignment $R^2$) for seat
    counts is $1$--$12\%$ under the soft POU; with hard nearest-word
    assignment it reaches only $\approx 0.26$--$0.50$ at $K=80$, where
    the 1{,}038 strata hold $\approx 17$ plans each.
  \item \textbf{Rare events are localized.} Six 2008 races have
    genuine seat variation; four races are rare events (a Democratic
    seat exists in $0.1$--$1\%$ of plans). At $K=12$ the top-3 strata
    contain $91$--$100\%$ of the rare-event plans, even though
    per-stratum enrichment is modest. All 2010--2016 races are flat
    (0 Democratic seats in every plan) --- no line-drawing choice can
    change them.
\end{itemize}

These numbers frame both projects: there is a lot of alignment $R^2$
``on the table'' for Project~1, and Project~2 has exact rare-event
probabilities to benchmark against.

\section{Project 1: What does the stratification see?\\
\normalsize Alignment of geometric strata with election outcomes}

\subsection*{Question}

The pipeline's strata are built from geometry alone. Are they
informative about the observables ensembles are actually used for ---
seat counts, partisan measures~\cite{stephanopoulos2015,katz2020},
majority-minority structure? If not, can election-aware district
metrics produce better-aligned strata without destroying spatial
interpretability?

\subsection*{Why the question matters}

Stratified sampling reduces the variance of $\hat{\mathbb{E}}_\pi[f]$
only through the between-stratum component of $\mathrm{Var}_\pi(f)$.
If the strata are nearly independent of $f$, stratification cannot
help, no matter how well the sampler is engineered. The alignment
$R^2$ is therefore the design quantity for the whole program, and the
Duplin--Onslow numbers above show it is small for the natural
geometric metric. Whether that is intrinsic (seat outcomes depend on
fine boundary details that any coarse geometric summary discards) or
fixable (a vote-aware metric recovers it) is open, and either answer
matters.

\subsection*{Stage 1 --- exact decomposition on Duplin--Onslow}

For each observable $f$ (six competitive 2008 seat counts, four
rare-event indicators, max district BVAP share), decompose
\[
  \mathrm{Var}_\pi(f)
  \;=\;
  \underbrace{\textstyle\sum_s z_s\,(\mathbb{E}_{\pi_s}[f]-\mathbb{E}_\pi[f])^2}_{\text{between}}
  \;+\;
  \underbrace{\textstyle\sum_s z_s\,\mathrm{Var}_{\pi_s}(f)}_{\text{within}},
  \qquad
  \frac{d\pi_s}{d\pi}=\frac{\phi_s}{z_s},
\]
exactly, sweeping the POU temperature $T$, beam width $L$, and
alphabet size $K$. Reproduce and extend the baseline table in
\texttt{local/duplin/alignment\_R2\_K12\_L10.csv}. Characterize the
$R^2$-vs-(number of strata) tradeoff curve and the rare-event
concentration $\max_s P(\text{event}\mid s)$. This stage is mostly
running and interrogating existing code; first plots within two weeks.

\subsection*{Stage 2 --- election-aware letters}

Augment the district distance with a vote-share term,
\[
  d_\lambda(u,v) \;=\; d_{\mathrm{geom}}(u,v)
  \;+\; \lambda\, W \,\bigl|\,\mathrm{share}(u)-\mathrm{share}(v)\,\bigr|,
\]
re-run \hccultra{} and the word construction across $\lambda$ (minutes
of compute at this scale), and trace alignment-vs-$\lambda$ for each
observable. The tension to characterize: as $\lambda$ grows, do
letters remain spatially interpretable district types, or fragment
into partisan bins? Quantify interpretability (e.g., spatial
coherence of letter densities) so the tradeoff curve has two real
axes.

\subsection*{Stage 3 --- does it transfer? (full scale)}

Repeat Stage 1 on the Connecticut ensemble from the paper (50K Cycle
Walk samples, $K=34$, statewide vote overlays from VEST), where the
decomposition is no longer exact but the sample is large. Compare: is
the weak-alignment finding a quirk of the lopsided Duplin--Onslow
electorate (all 2010--2016 races are flat there), or does it persist
in a competitive state? This stage turns the exact small-scale finding
into a paper-ready claim.

\subsection*{Deliverables and skills}

A reproducible alignment analysis (notebook + figures) for both
scales; an alignment section that slots into the follow-up paper.
Requires Python/pandas and basic probability (law of total variance);
no MCMC internals. The pipeline is used as a black box.

\section{Project 2: Does stratified sampling work?\\
\normalsize Implementing and benchmarking the umbrella sampler}

\subsection*{Question}

Implement the stratified sampler the paper sets up --- tilted chains
per stratum, EMUS recombination~\cite{thiede2016,dinner2020} --- and
benchmark it where the right answer is known exactly.

\subsection*{Why Duplin--Onslow first}

A sampler benchmark is only as good as its ground truth. Here the
exact stratum weights $z_s$, the exact flux matrix, and the exact
rare-event probabilities (e.g.\ $P(\text{EL08G\_PR Dem seat}) \approx
0.002$) are already computed. Every component of the implementation
can be unit-tested against truth before any scaling question arises.
The $I=3$ district matching is also trivial at this scale, so the
sampler logic can be developed where the bookkeeping is easy.

\subsection*{Stage 1 --- pin down the tilt}

The draft writes $J_s = J - \phi_s$, which is at most an $e$-fold
reweighting (since $\phi_s\in[0,1]$) --- too weak to confine a chain
to a stratum. The standard umbrella choice is
$J_s = J - \log\psi(x;s)$, i.e.\ $\pi_s \propto \psi(\cdot\,;s)\,\pi$,
with unnormalized $\psi$ avoiding the POU normalizer in the
Metropolis--Hastings ratio. Resolving this (with the group) is the
first milestone and a genuine contribution to the writeup.

\subsection*{Stage 2 --- biased chains on Duplin--Onslow}

Modify the Cycle Walk acceptance step to include
$\psi(x';s)/\psi(x;s)$. Evaluating $\psi$ requires the plan-to-word
distance --- a min-cost bipartite matching~\cite{burkard2009} that can
be warm-started, since each move changes few districts. Run one chain
per stratum, estimate the overlap matrix $\hat F$, solve
$z^\top F = z^\top$, and attach EMUS asymptotic error
bars~\cite{thiede2016}. \textbf{Validation gates:} $\hat z_s \to z_s$
(exact), $\hat F \to F$ (exact), estimated observables $\to$ exact
values, with error bars that cover at the nominal rate.

\subsection*{Stage 3 --- the benchmark}

Headline plot: MSE vs.\ total compute for the stratified estimator
against a single long Cycle Walk chain, for (a) a smooth observable
(max BVAP share, or a competitive 2008 seat count) and (b) a
rare-event probability, with the exact answer drawn as a horizontal
line. The rare events are where the win should be dramatic: the top-3
strata contain 91--100\% of the rare-event plans, so chains tilted
into those strata see the event constantly while the unstratified
chain waits for a $10^{-2}$--$10^{-3}$ fluctuation.

\subsection*{Stage 4 --- climb the ladder (full scale)}

Move up the grid-graph sequence ($6\times6$ through $20\times20$,
pre-generated in \texttt{data/graph/}) to where enumeration is
impossible and stratification has to earn its keep, tracking how the
matching cost, overlap structure, and variance reduction scale with
$|V|$ and $I$. Connecticut ($I=5$) is the graduation exercise and the
bridge back to the paper; the incremental-matching machinery developed
here is what makes $I=14$ (North Carolina) feasible later.

\subsection*{Deliverables and skills}

A working, validated stratified sampler in the \texttt{gmstrat}
repository; the MSE-vs-compute benchmark at both scales; a methods
section for the follow-up paper. Requires stronger programming
(Python plus some Julia for the Cycle Walk harness) and comfort with
MCMC. Main schedule risk is ramping up on the Cycle Walk codebase ---
pair with Zijian early.

\section{How the projects interlock}

Project 1 tells Project 2 which observables stratification can
plausibly help with: if geometric strata cannot see seat counts,
Project 2 benchmarks against Project 1's election-aware strata, and
the comparison itself is the joint result. Both run on infrastructure
that exists today; neither blocks the other. The natural merge point
is a follow-up paper --- stratified sampling for election-outcome
estimation, validated exactly at small scale --- with both students on
it.

\section{Infrastructure inventory}\label{sec:infra}

All in the \texttt{gmstrat} repository (branch
\texttt{grid-experiments}):

\begin{itemize}[itemsep=2pt]
  \item \texttt{duplin.py} --- Duplin--Onslow converters, a fast
    \hccultra{} implementation (validated to produce byte-identical
    linkage against \texttt{hccfit.py}; full 14{,}345-district run in
    $\approx 3$ minutes), election utilities.
  \item \texttt{data/duplin\_onslow/} --- raw Duke data, gmstrat-format
    graph, all 17{,}653 plans in Cycle Walk JSONL format.
  \item \texttt{local/duplin/} --- cached pipeline artifacts (distance
    matrix, linkage, per-plan seat counts), the elbow and
    stratification drivers (\texttt{run\_elbow.py},
    \texttt{run\_strata.py}), figures, and the baseline alignment
    table.
  \item \texttt{data/graph/} --- the grid ladder; \texttt{sampling/}
    --- the Julia Cycle Walk harness; \texttt{demo.ipynb} --- the
    canonical end-to-end pipeline walkthrough.
  \item The annotated reading list (with per-project reading orders)
    lives in the Overleaf project under
    \texttt{undergrad\_reading/}.
\end{itemize}

\begin{thebibliography}{9}

\bibitem{cyclewalk}
D.~R. DeFord, G.~Herschlag, J.~C. Mattingly,
\emph{A cycle walk for sampling measures on spanning forests for
redistricting}, arXiv:2509.08629, 2025.

\bibitem{yim2024}
J.-H. Yim, A.~C. Gilbert,
\emph{Fitting trees to $\ell_1$-hyperbolic distances},
arXiv:2409.01010, 2024.

\bibitem{herschlag2020}
G.~Herschlag, H.~S. Kang, J.~Luo, C.~V. Graves, S.~Bangia,
R.~Ravier, J.~C. Mattingly,
\emph{Quantifying gerrymandering in North Carolina},
Statistics and Public Policy 7(1):30--38, 2020.

\bibitem{thiede2016}
E.~H. Thiede, B.~Van Koten, J.~Weare, A.~R. Dinner,
\emph{Eigenvector method for umbrella sampling enables error
analysis}, J.~Chem.\ Phys.\ 145, 084115, 2016.

\bibitem{dinner2020}
A.~R. Dinner, E.~H. Thiede, B.~Van Koten, J.~Weare,
\emph{Stratification as a general variance reduction method for
Markov chain Monte Carlo}, SIAM/ASA J.~Uncertainty Quantification
8(3):1139--1188, 2020.

\bibitem{stephanopoulos2015}
N.~O. Stephanopoulos, E.~M. McGhee,
\emph{Partisan gerrymandering and the efficiency gap},
University of Chicago Law Review 82(2):831--900, 2015.

\bibitem{katz2020}
J.~N. Katz, G.~King, E.~Rosenblatt,
\emph{Theoretical foundations and empirical evaluations of partisan
fairness in district-based democracies},
American Political Science Review 114(1):164--178, 2020.

\bibitem{burkard2009}
R.~Burkard, M.~Dell'Amico, S.~Martello,
\emph{Assignment Problems}, SIAM, 2009.

\end{thebibliography}

\end{document}
